\documentclass[12pt, a4paper]{article}
\usepackage{geometry}
\usepackage{amsmath}
\usepackage{proof}
\usepackage{amssymb}
\usepackage{enumitem}
\geometry{a4paper, margin=1in}

\title{Trabajo práctico 1 \\ \large Análisis de Lenguajes de Programación}

\author{Sofía Grepachok (G-5949/8) \and Thelma Medrán (M-7412/8)}
\date{Septiembre 2026}

\begin{document}

\maketitle

%%%%%%%%%%%%%%%%%%%%%%%%%%%%%%%%%%%%%%%%%%%%%%%%%%%%%%%%%%%%%%%%%%%%%%%%%%%%%%%%

\section*{Ejercicio 1}

Extendemos la sintaxis abstracta de LIS de la siguiente manera:
\[
\begin{array}{rcl}
\mathit{intexp} & ::= & \mathit{nat} \\
                & \mid & \mathit{var} \\
                & \mid & -_u \mathit{intexp} \\
                & \mid & \mathit{intexp} + \mathit{intexp} \\
                & \mid & \mathit{intexp} -_b \mathit{intexp} \\
                & \mid & \mathit{intexp} \times \mathit{intexp} \\
                & \mid & \mathit{intexp} \div \mathit{intexp} \\
                & \mid & \mathit{var} \mathbf{++} \\
                & \mid & \mathit{var} \mathbf{--} 
\end{array}
\]
Análogamente, extendemos la sintaxis concreta:
\[
\begin{array}{rcl}
\mathit{intexp} & ::= & \mathit{nat} \\
                & \mid & \mathit{var} \\
                & \mid & \text{'\(-\)' } \mathit{intexp} \\
                & \mid & \mathit{intexp} \text{ '\(+\)' } \mathit{intexp} \\
                & \mid & \mathit{intexp} \text{ '\(-\)' } \mathit{intexp} \\
                & \mid & \mathit{intexp} \text{ '\(*\)' } \mathit{intexp} \\
                & \mid & \mathit{intexp} \text{ '\(/\)' } \mathit{intexp} \\
                & \mid & \text{ '\((\)' } \mathit{intexp} \text{ '\()\)' } \\
                & \mid & \mathit{var} \text{ '\(++\)'} \\
                & \mid & \mathit{var} \text{ '\(--\)'}                 
\end{array}
\]
Por último, la nueva lista de precedencia queda como sigue:

\begin{center}
\begin{minipage}{0.3\textwidth}
\begin{enumerate}
\setlength\itemsep{0pt}
\setlength\parskip{0pt}
    \item $(\mathbf{++} \quad \mathbf{--} \quad -_u)$
    \item $(* \quad /)$
    \item $(+ \quad -_b)$
    \item $(== \quad \neq \quad < \quad >)$
    \item $!$
    \item $\&\&$
    \item $\mid \mid$
    \item $=$
    \item $;$
\end{enumerate}
\end{minipage}
\end{center}

%%%%%%%%%%%%%%%%%%%%%%%%%%%%%%%%%%%%%%%%%%%%%%%%%%%%%%%%%%%%%%%%%%%%%%%%%%%%%%%%
\section*{Ejercicio 3}

Para parsear expresiones enteras debemos primero desambiguar la grámatica para respetar el orden de precendencia y asociar a izquierda: 
\[
\begin{array}{lcl}
exp  & ::= & exp\ \texttt{'+'}\ term \mid exp\ \texttt{'$-_b$'}\ term \mid term \\
term & ::= & term\ \texttt{'*'}\ unary \mid term\ \texttt{'/'}\ unary \mid unary \\
unary & ::= & \ \texttt{'$-_u$'}\ unary \mid atom \\
atom & ::= & \texttt{'('}\ exp\ \texttt{')'} \mid nat \mid var\ \texttt{'++'} \mid var\ \texttt{'--'} \\
digit  & ::= & \texttt{'0'} \mid \texttt{'1'} \mid \cdots \mid \texttt{'9'} \\
letter & ::= & \texttt{'a'} \mid \cdots \mid \texttt{'Z'} \\
nat  & ::= & digit \mid digit\ nat \\
var  & ::= & letter \mid letter\ var \\
\end{array}
\]

Definimos el parser con respecto a esto.

\section*{Ejercicio 4}

Agregamos las extensiones del Ejercicio 1 a la semántica big-step:

\[
\infer[\textsc{VarInc}]
  {\langle x{+}{+}, \sigma\rangle \Downarrow_{exp} \langle n+1, [\sigma \mid x : n+1]\rangle}
  {\langle x, \sigma\rangle \Downarrow_{exp} \langle n, \sigma\rangle}
\]

\[
\infer[\textsc{VarDec}]
  {\langle x{-}{-}, \sigma\rangle \Downarrow_{exp} \langle n-1, [\sigma \mid x : n-1]\rangle}
  {\langle x, \sigma\rangle \Downarrow_{exp} \langle n, \sigma\rangle}
\]

%%%%%%%%%%%%%%%%%%%%%%%%%%%%%%%%%%%%%%%%%%%%%%%%%%%%%%%%%%%%%%%%%%%%%%%%%%%%%%%%

\section*{Ejercicio 5}


%%%%%%%%%%%%%%%%%%%%%%%%%%%%%%%%%%%%%%%%%%%%%%%%%%%%%%%%%%%%%%%%%%%%%%%%%%%%%%%%

\section*{Ejercicio 6}

Por simplicidad, definimos
\[
c_2 : y\ {=}\ x{+}{+}
\qquad
c_3 : \mathbf{if}\ x\ {==}\ 0\ \mathbf{then}\ z\ {=}\ 1\ \mathbf{else}\ z\ {=}\ 0
\]

Mostramos que $\langle x\ {=}\ 0;\ c_2;\ c_3,\ \{\}\rangle \leadsto^* \langle \mathbf{skip},\ \{x\ {:}\ 1,\ y\ {:}\ 1,\ z\ {:}\ 0\}\rangle$ en seis pasos.

\begin{enumerate}[leftmargin=*]

\item Primero, probamos que $\langle x\ {=}\ 0;\ c_2;\ c_3,\ \{\}\rangle \leadsto \langle \mathbf{skip};\ c_2;\ c_3,\ \{x\ {:}\ 0\}\rangle$:

\[
\infer[\textsc{Seq2}]
  {\langle x\ {=}\ 0;\ c_2;\ c_3,\ \{\}\rangle \leadsto \langle \mathbf{skip};\ c_2;\ c_3,\ \{x\ {:}\ 0\}\rangle}
  {\infer[\textsc{Ass}]
    {\langle x\ {=}\ 0,\ \{\}\rangle \leadsto \langle \mathbf{skip},\ \{x\ {:}\ 0\}\rangle}
    {\infer[\textsc{NVal}]
      {\langle 0,\ \{\}\rangle \Downarrow_{exp} \langle 0,\ \{\}\rangle}
      {}}}
\]

\item Luego, probamos que $\langle \mathbf{skip};\ c_2;\ c_3,\ \{x\ {:}\ 0\}\rangle \leadsto \langle c_2;\ c_3,\ \{x\ {:}\ 0\}\rangle$:

\[
\infer[\textsc{Seq1}]
  {\langle \mathbf{skip};\ c_2;\ c_3,\ \{x\ {:}\ 0\}\rangle \leadsto \langle c_2;\ c_3,\ \{x\ {:}\ 0\}\rangle}
  {}
\]

\item Ahora, veamos que $\langle y\ {=}\ x{+}{+};\ c_3,\ \{x\ {:}\ 0\}\rangle \leadsto \langle \mathbf{skip};\ c_3,\ \{x\ {:}\ 1,\ y\ {:}\ 1\}\rangle$:

\[
\infer[\textsc{Seq2}]
  {\langle y\ {=}\ x{+}{+};\ c_3,\ \{x\ {:}\ 0\}\rangle \leadsto \langle \mathbf{skip};\ c_3,\ \{x\ {:}\ 1,\ y\ {:}\ 1\}\rangle}
  {\infer[\textsc{Ass}]
    {\langle y\ {=}\ x{+}{+},\ \{x\ {:}\ 0\}\rangle \leadsto \langle \mathbf{skip},\ \{x\ {:}\ 1,\ y\ {:}\ 1\}\rangle}
    {\infer[\textsc{VarInc}]
      {\langle x{+}{+},\ \{x\ {:}\ 0\}\rangle \Downarrow_{exp} \langle 1,\ \{x\ {:}\ 1\}\rangle}
      {\infer[\textsc{Var}]
        {\langle x,\ \{x\ {:}\ 0\}\rangle \Downarrow_{exp} \langle 0,\ \{x\ {:}\ 0\}\rangle}
        {x \in \mathrm{dom}\,\{x\ {:}\ 0\}}}}}
\]

\item Y se tiene que $\langle \mathbf{skip};\ c_3,\ \{x\ {:}\ 1,\ y\ {:}\ 1\}\rangle \leadsto \langle c_3,\ \{x\ {:}\ 1,\ y\ {:}\ 1\}\rangle$:

\[
\infer[\textsc{Seq1}]
  {\langle \mathbf{skip};\ c_3,\ \{x\ {:}\ 1,\ y\ {:}\ 1\}\rangle \leadsto \langle c_3,\ \{x\ {:}\ 1,\ y\ {:}\ 1\}\rangle}
  {}
\]

\item Probamos que $\langle \mathbf{if}\ x\ {==}\ 0\ \mathbf{then}\ \{z\ {=}\ 1\}\ \mathbf{else}\ \{z\ {=}\ 0\},\ \{x\ {:}\ 1,\ y\ {:}\ 1\}\rangle \leadsto \langle z{=}0,\ \{x\ {:}\ 1,\ y\ {:}\ 1\}\rangle$:

\begin{center}
\infer[\textsc{If2}]
  {\langle \mathbf{if}\ x\ {==}\ 0\ \mathbf{then}\ \{z\ {=}\ 1\}\ \mathbf{else}\ \{z\ {=}\ 0\},\ \{x\ {:}\ 1,\ y\ {:}\ 1\}\rangle \leadsto \langle z\ {=}\ 0,\ \{x\ {:}\ 1,\ y\ {:}\ 1\}\rangle}
  {\infer[\textsc{Eq}]
    {\langle x\ {==}\ 0,\ \{x\ {:}\ 1,\ y\ {:}\ 1\}\rangle \Downarrow_{exp} \langle \mathit{false},\ \{x\ {:}\ 1,\ y\ {:}\ 1\}\rangle}
    {\infer[\textsc{Var}]
      {\langle x,\ \{x\ {:}\ 1,\ y\ {:}\ 1\}\rangle \Downarrow_{exp} \langle 1,\ \{x\ {:}\ 1,\ y\ {:}\ 1\}\rangle}
      {x \in \mathrm{dom}\,\{x\ {:}\ 1,\ y\ {:}\ 1\}}
     &
     \infer[\textsc{NVal}]
       {\langle 0,\ \{x\ {:}\ 1,\ y\ {:}\ 1\}\rangle \Downarrow_{exp} \langle 0,\ \{x\ {:}\ 1,\ y\ {:}\ 1\}\rangle}
       {}}}
\end{center}
\item Y por último, probamos que $\langle z\ {=}\ 0,\ \{x\ {:}\ 1,\ y\ {:}\ 1\}\rangle \leadsto \langle \mathbf{skip},\ \{x\ {:}\ 1,\ y\ {:}\ 1,\ z\ {:}\ 0\}\rangle$:

\[
\infer[\textsc{Ass}]
  {\langle z\ {=}\ 0,\ \{x\ {:}\ 1,\ y\ {:}\ 1\}\rangle \leadsto \langle \mathbf{skip},\ \{x\ {:}\ 1,\ y\ {:}\ 1,\ z\ {:}\ 0\}\rangle}
  {\infer[\textsc{NVal}]
    {\langle 0,\ \{x\ {:}\ 1,\ y\ {:}\ 1\}\rangle \Downarrow_{exp} \langle 0,\ \{x\ {:}\ 1,\ y\ {:}\ 1\}\rangle}
    {}}
\]

\end{enumerate}

Encadenando los seis pasos anteriores, concluimos
\[
\langle x\ {=}\ 0;\ c_2;\ c_3,\ \{\}\rangle
\leadsto^*
\langle \mathbf{skip},\ \{x\ {:}\ 1,\ y\ {:}\ 1,\ z\ {:}\ 0\}\rangle
\]
\end{document}
